\section{Retained model objects}

We use the notation of Ishikawa and Shibata (2021). Let $a$ and $c$ denote the demand and baseline-cost parameters, $\gamma$ the R\&D-cost parameter, $\lambda$ the relative-profit/market-competitiveness parameter, and $\beta_1,\beta_2$ the asymmetric spillover parameters. Denote by
\[
 x_i^n=x_i^n(a,c,\gamma,\lambda,\beta_1,\beta_2)
\]
the published noncooperative R\&D equilibrium of firm $i$, and by
\[
 x_i^c=x_i^c(a,c,\gamma,\lambda,\beta_1,\beta_2)
\]
the corresponding cooperative R\&D equilibrium. These are the closed-form solutions reported in Eqs. (16)--(20) and (25)--(29) of the original paper. We retain those solutions rather than re-derive the full model here.

For the individual comparison define
\begin{equation}\label{eq:Di}
 \Di(\beta_1,\beta_2,\lambda,\gamma)=x_i^n-x_i^c.
\end{equation}
An individual equality boundary is therefore a zero set $\Di=0$ on a regular region in which the retained closed forms are defined. For the aggregate comparison define
\begin{equation}\label{eq:DA}
 \DA(\beta_1,\beta_2,\lambda,\gamma)
 =(x_1^n+x_2^n)-(x_1^c+x_2^c).
\end{equation}
The aggregate equality boundary is $\DA=0$. The signs of these objects compare R\&D levels; they are not, by themselves, social-welfare rankings. All results below are restricted to regular parameter regions in which the denominators of the retained closed forms are nonzero and the stated interior solutions satisfy the maintained admissibility conditions. The sign convention is straightforward: $D_i>0$ means that firm $i$ undertakes more R\&D under noncooperation than under cooperation, while $D_A>0$ means that aggregate R\&D is larger under noncooperation. The equality sets are the boundaries between these sign regions. A claim that a threshold is independent of $\gamma$ is therefore a claim about the geometry of an entire zero set, not merely about the value of $D_i$ or $D_A$ at one calibration. Likewise, a claim that the aggregate boundary depends only on average spillovers asserts that moving along a constant-sum line cannot change the equality condition. Propositions 1 and 2 test precisely these geometric implications.

This formulation isolates the source of the correction. Nothing in Propositions 1--3 requires a different equilibrium concept or a modification of the original game. The question is whether the zero sets generated by the retained equilibrium formulas have the parameter-invariance and scalar-threshold properties attributed to them in the subsequent numerical analysis. Algebraically, this means that the numerators obtained after clearing nonzero denominators must be treated as equations in the relevant parameters, rather than summarized by a calibration-specific crossing. Economically, it means that a comparison between competition and cooperation should be indexed by the actual zero set of the retained R\&D choices. The regular-domain qualification is therefore substantive: clearing a denominator does not authorize conclusions at a singular point, and every local argument below is paired with a nonzero-denominator certificate.
